\section{Method}


\subsection{Speculative Reasoning}
\label{sec:spec_reason}

Due to its reliance on autoregressive decoding, LRM inference incurs significantly higher latency than typical LLMs---often to the point of being prohibitively slow for interactive applications and degrading user experience~\citep{dynasor_cot}. Existing approaches for latency reduction include using a distilled version of the base model~\citep{deepseek_r1}, limiting the number of thinking tokens via a predefined {\em token budget}, or disabling the reasoning process altogether by omitting the thinking tokens (\texttt{<think>} and \texttt{</think>}) during generation~\citep{qwen3}. However, these approaches impose a harshly trade-off between accuracy for latency: they either limit the model's capacity to reason or apply a lower-quality model uniformly across all reasoning steps.
In contrast, \name{} takes a more fine-grained and adaptive approach. Instead of explicitly restricting output length, it selectively offloads only the easier reasoning steps to a lightweight model, preserving overall reasoning quality while substantially reducing inference latency.

The approximation-tolerant nature of LRM reasoning enables a new form of speculative execution: tentatively carrying out reasoning steps using a lightweight model, assessing their utility with a stronger base model, and selectively accepting them. \name{} leverages this flexibility to reduce decoding latency while preserving output quality.
To achieve this goal, \name{} offloads easier or less critical reasoning steps---defined as semantically self-contained units such as complete sentences or logical steps---to a smaller, faster {\em speculator} model. Each step is decoded in two stages: (1) the lightweight speculator proposes the next reasoning step based on the current context, and (2) the base model evaluates the proposed step for semantic utility. If the step is accepted, \name{} proceeds to the next step; otherwise, \name{} falls back to the base model to regenerate the step. While our implementation uses a simple static-threshold mechanism for verification, the framework supports richer, customizable decision strategies. We outline key design principles below.

\textbf{Navigating the Pareto frontier of the latency-accuracy tradeoff.} \name{} expands the Pareto frontier of the latency-accuracy tradeoff by exposing fine-grained control knobs to navigate through this space. The key knob \name{} employs is the acceptance threshold: after each speculated reasoning step, the base model is prompted to generate a single-token utility score (e.g., an integer from 0 to 9) indicating the quality of the step. If the utility score is above a static acceptance threshold (e.g., score $\geq 7$), the speculated reasoning step is accepted; otherwise, it is discarded and regenerated by the base model.

Adjusting this threshold allows users to control the {\em strictness} of speculation (Fig.~\ref{fig:lat_acc_tradeoff}): a higher threshold requires speculated steps to be closer to token-level equivalence on the equivalence spectrum (Fig.~\ref{fig:equiv_spectrum}), improving accuracy but reducing the acceptance rate and thereby increasing latency. Conversely, a lower threshold increases speculation efficiency at the cost of potential accuracy degradation.

An additional knob involves forcing the first $n$ reasoning steps to be decoded by the base model. Since LRMs often use the initial steps to analyze the problem and formulate a high-level plan, assigning these initial steps to the base model can steer the overall reasoning trajectory toward higher quality.
We show in Fig.~\ref{fig:earlyforce_knob} that this knob also allows \name{} to manage latency-accuracy tradeoff, though with less impact than the acceptance threshold knob.


While our current implementation uses a simple, discrete threshold-based scoring scheme---offering only a coarse-grained configuration space---it establishes a lower bound on verification quality. 
Future work can explore more sophisticated strategies, such as logprob-based confidence estimates or dynamic thresholds, to enable finer-grained tradeoffs without incurring additional runtime cost, and may further improve overall performance.










\textbf{Efficient verification.} Because each step requires verification by the base model, it's crucial to keep verification overhead low to avoid compounding latency. 
Instead of autoregressively decoding or reranking multiple candidate steps, \name{} evaluates each speculated step in a single \textit{prefill-only} pass of the base model. The verification prompt is templated to reuse most of the CoT prefix, so each verification requires prefilling only $\sim$70 new tokens. Since short-prefill forward passes are memory-bound, the overhead is comparable to decoding just 1–2 tokens, making verification highly efficient in practice.

\textbf{Implementation details.} Since the small model is lightweight, we colocate both the small and base models on the same GPU. The memory reserved for Key-Value caches~\citep{vllm} is statically partitioned between the two models. They do not share any internal model states--only the token IDs of the generated reasoning steps are managed and shared by \name{}. If a speculative step is rejected, the corresponding KV cache entries are discarded. 

Inference is performed sequentially: the small and base models take turns, avoiding kernel-level interference. In future work, we plan to explore pipelining to overlap the small model’s decoding with the base model’s inference. While this may introduce mild resource contention, it could further reduce end-to-end latency.








\subsection{Hierarchical Speculation across Semantic Similarity and Token Equivalence}
\label{sec:hierarchical_speculation}


At a high level, \name{}'s speculative reasoning resembles the philosophy behind traditional speculative decoding, but differs in two important ways. First, speculative decoding guarantees token-level equivalence between draft and verified outputs, making it a form of exact acceleration. In contrast, \name{} targets semantic-level similarity, accepting steps that carry the same insight even if phrased differently, and exposes knobs to control the exactness of reasoning approximations. Second, speculative decoding is typically applied to output generation tasks (e.g., text continuation or translation), where the fidelity of each token matters. \name{}, on the other hand, is designed specifically for internal thinking tokens in reasoning tasks, where intermediate steps are approximate and interchangeable as long as they preserve the logical progression of thought.




Further, because \name{} and speculative decoding operate at different levels (semantic-level similarity vs. token-level equivalence), these two approaches are complementary and can be combined into a unified, hierarchical system -- \name{}+Decode first applies step-level speculative reasoning to draft and verify reasoning steps. If a step is rejected and regenerated by the base model, standard token-level speculative decoding can be applied during the base model regeneration to further accelerate decoding. %







